\setcounter{section}{0}

\Needspace{5\baselineskip}
\hypertarget{ux952eux503cux7f13ux5b58kv-cache}{%
\section{键值缓存（KV Cache）}\label{ux952eux503cux7f13ux5b58kv-cache}}

\Needspace{5\baselineskip}
\hypertarget{ux4e00kv-cacheux7684ux6838ux5fc3ux539fux7406}{%
\subsection{一、KV Cache的核心原理}\label{ux4e00kv-cacheux7684ux6838ux5fc3ux539fux7406}}

在全注意力（Full Attention）中，生成每一个token的时间、空间复杂度为\(O(n^2)\)（其中在推理阶段运用KV Cache能让生成单个token的时间复杂度降为\(O(n)\)、生成序列的时间复杂度为\(O(n^2)\)，内存空间复杂度为\(O(n)\)），这些成为处理长序列（长上下文）的主要障碍。特别是在推理阶段（不是训练阶段），生成每一个token，计算的成本均为\(O(L^2)\)。然而在通过输入序列计算K、V矩阵时，这里很多都是重复计算。

\Needspace{5\baselineskip}
为了解决推理时计算时间复杂度的问题，LLM普遍运用了KV Cache：

\Needspace{4\baselineskip}
\begin{enumerate}
\def\labelenumi{\arabic{enumi}.}
\tightlist
\item
  由输入计算Q、K、V时
\end{enumerate}

\Needspace{5\baselineskip}
传统注意力机制中：

\[
Q = XW_Q,\quad K = XW_K,\quad V = XW_V
\]

\Needspace{5\baselineskip}
而KV Cache中，由于在生成第t+1个token时，我们只需要第t个token作为Q。而对于K和V，其实在之前生成中已经计算过其第1至t-1行（即第1至t-1个token的K和V），可以把它们存在内存中，从而可以直接拿来用，只需要计算第t个token的K和V：

\[
q_t = x_tW_Q,\quad k_t = x_tW_K,\quad v_t = x_tW_V
\]

这里的k\_t和v\_t在内存里拼接到之前算出的K和V矩阵下面，供后续使用。这样就能将计算复杂度从O(L)变成O(1)。

\Needspace{4\baselineskip}
\begin{enumerate}
\def\labelenumi{\arabic{enumi}.}
\setcounter{enumi}{1}
\tightlist
\item
  计算注意力权重时
\end{enumerate}

\Needspace{5\baselineskip}
在传统注意力机制中：

\[
\mathrm{Scores} = Q \cdot K^T
\]

\Needspace{5\baselineskip}
操作描述：

\begin{itemize}
\tightlist
\item
  \(Q\) 的维度是 \(t \times d_k\)。
\item
  \(K\) 的维度是 \(t \times d_k\)，因此 \(K^T\) 的维度是 \(d_k \times t\)。
\item
  这是一个标准的矩阵乘法，得到的注意力分数矩阵维度是 \(t \times t\)。
\end{itemize}

\Needspace{5\baselineskip}
在KV Cache中：

\[
\mathrm{Scores} = q_t \cdot (K_{\mathrm{cache}}^{(t)})^T
\]

\Needspace{5\baselineskip}
操作描述：

\begin{itemize}
\tightlist
\item
  \(q_t\) 的维度是 \(1 \times d_k\)。
\item
  \(K_{\mathrm{cache}}^{(t)}\) 的维度是 \(t \times d_k\)（转置后为 \(d_k \times t\)）。
\item
  这是一次``向量-矩阵''乘法，得到的当前步注意力分数维度是 \(1 \times t\)。
\end{itemize}

从而计算复杂度从O(L\^{}2)变成O(L)

\Needspace{4\baselineskip}
\begin{enumerate}
\def\labelenumi{\arabic{enumi}.}
\setcounter{enumi}{2}
\tightlist
\item
  加权得到输出值时
\end{enumerate}

\Needspace{5\baselineskip}
在传统注意力中：

\[
Y = \mathrm{Softmax}(\mathrm{Scores}) \cdot V
\]

\Needspace{5\baselineskip}
操作描述：

\begin{itemize}
\tightlist
\item
  注意力矩阵维度是 \(t \times t\)。
\item
  \(V\) 矩阵维度是 \(t \times d_v\)。
\item
  输出 \(Y\) 的维度是 \(t \times d_v\)。
\end{itemize}

\Needspace{5\baselineskip}
在KV Cache中：

\[
y_t = \alpha_t \cdot V_{\mathrm{cache}}^{(t)}
\]

\Needspace{5\baselineskip}
操作描述：

\begin{itemize}
\tightlist
\item
  权重 \(\alpha_t\) 的维度是 \(1 \times t\)。
\item
  \(V_{\mathrm{cache}}^{(t)}\) 的维度是 \(t \times d_v\)。
\item
  这是一次``向量-矩阵''乘法，输出 \(y_t\) 的维度是 \(1 \times d_v\)。
\end{itemize}

\Needspace{5\baselineskip}
\hypertarget{ux4e8cux4e3aux4ec0ux4e48ux8badux7ec3ux7684ux65f6ux5019ux4e0dux7528kv-cache}{%
\subsection{二、为什么训练的时候不用KV Cache？}\label{ux4e8cux4e3aux4ec0ux4e48ux8badux7ec3ux7684ux65f6ux5019ux4e0dux7528kv-cache}}

KV Cache是为了解决推理（Inference）阶段每次生成一个token的``串行计算''低效问题而设计的；而训练（Training）阶段天然是并行计算的，我们已经知道完整答案，使用Teacher Forcing和掩码技术，可以在一次计算中算出所有位置的Loss，自然就没有``缓存上一步结果给下一步用''的需求。如果在训练阶段强行使用KV Cache，反而会把高效的并行训练退化为低效的串行训练，导致训练速度下降成百上千倍。

前面介绍过，预训练或微调阶段，各个步骤是在一个矩阵运算中同时完成的。

我们一次性拿到了所有token 的 Q、K、V，随后计算注意力分数，并进行掩码。掩码是一个``上三角矩阵''，它在计算注意力分数后、进行 Softmax 之前起作用。它会强制将所有未来位置(i,j)（i\textgreater j）的注意力分数，即q\_i与k\_j的点积设置为一个绝对值极大的负数（例如 -1e9）。

这个过程中不存在``第1步算完存下来给第2步用''的情况，因为第1步和第2步是同时发生的。

【特例：在RLHF的PPO训练阶段，流程分为两步：Rollout（生成阶段）让模型针对 Prompt 生成回答，这一步本质是推理，所以这时候会用到 KV Cache；Update（更新阶段）拿着生成的回答和奖励（Reward）去计算梯度更新模型，回到并行模式，不用KV Cache。】

\Needspace{5\baselineskip}
\hypertarget{ux4e09ux9884ux586bux5145prefill}{%
\subsection{三、预填充（Prefill）}\label{ux4e09ux9884ux586bux5145prefill}}

在生成开始前，我们会让模型一次性并行读完所有prompt（提示词），利用Q = X\emph{W\_Q，K = X}W\_K，V=X*W\_V计算初始Q、K、V矩阵，将初始状态下的KV Cache存入缓存。然后，再开始解码（Decode），自回归地生成token。

预填充是计算密集型任务，而解码是内存密集型任务。在一些架构中，会让二者分别用高算力和高显存带宽的GPU。

\Needspace{5\baselineskip}
\hypertarget{ux56dbux7cfbux7edfux786cux4ef6ux7ba1ux7406paged-attentionvllm}{%
\subsection{四、系统硬件管理：Paged Attention（vLLM）}\label{ux56dbux7cfbux7edfux786cux4ef6ux7ba1ux7406paged-attentionvllm}}

KV Cache在显存中需要连续空间，导致大量的显存碎片浪费（类似于操作系统的内存碎片）。借鉴操作系统的虚拟内存页（Paging）技术。将KV Cache分块存储在不连续的物理显存中，通过页表来索引。这极大地提高了显存利用率，使得在同样的显存下可以支持更大的Batch Size，从而提升推理吞吐量。

\FloatBarrier
\Needspace{5\baselineskip}
\hypertarget{ux53c2ux8003ux6587ux732e-91}{%
\subsection{参考文献}\label{ux53c2ux8003ux6587ux732e-91}}

\vspace{0.25em}
\begin{itemize}
\tightlist
\item
  Kwon, W., Li, Z., Zhuang, S., et al.~(2023). \href{https://arxiv.org/abs/2309.06180}{Efficient Memory Management for Large Language Model Serving with PagedAttention}. SOSP 2023.
\end{itemize}

\clearpage
